\documentclass[../../../include/open-logic-section]{subfiles}

\begin{document}

\olfileid{sth}{replacement}{extrinsic}
\olsection[বাহ্যিক বিবেচনা]{প্রতিস্থাপন সম্পর্কে বাহ্যিক বিবেচনা}

প্রতিস্থাপনের ন্যায্যতা প্রতিষ্ঠার একটি
\emph{বাহ্যিক} চেষ্টা দিয়ে শুরু করি।
বুলস এমন একটি যুক্তি দিয়েছেন:
\begin{quote}
  [\ldots] প্রতিস্থাপনের স্বতঃসিদ্ধগুলি
  গ্রহণের কারণ বেশ সরল: এগুলির অনেক
  কাম্য ফল আছে, আর (দৃশ্যত) কোনো অকাম্য
  ফল নেই। পর্যায়ক্রমিক ধারণা সম্পর্কে
  উপপাদ্য ছাড়াও এই ফলগুলির মধ্যে আছে
  অসীম সংখ্যার সন্তোষজনক—যদিও আদর্শ
  নয়—একটি তত্ত্ব এবং সুভিত্তিক সম্পর্কের
  উপর আবেশী সংজ্ঞাকে ন্যায্যতা দেওয়া
  অত্যন্ত কাম্য একটি ফল।
  \citep[229]{Boolos1971}
\end{quote}
বুলসের বক্তব্যের সার হলো, প্রতিস্থাপনের
ফল দেখে তাকে ন্যায্য প্রতিপন্ন করা উচিত।
তিনি যে নির্দিষ্ট ফলগুলির কথা বলেছেন,
সেগুলিই গত কয়েক অধ্যায়ে আলোচনা করেছি।
প্রতিস্থাপন দিয়ে দেখেছি, ফন নয়মান
ক্রমসংখ্যাগুলি সুক্রমবিন্যাসের প্রকারের
ধারণার চমৎকার প্রতিনিধি (এটাই আমাদের
``অসীম সংখ্যার সন্তোষজনক, যদিও আদর্শ নয়,
তত্ত্ব'')। প্রতিস্থাপন দিয়ে $V_\alpha$-গুলি
সংজ্ঞায়িত করেছি, স্তরাঙ্ক প্রতিষ্ঠা করেছি
এবং $\in$-আবেশ প্রমাণ করেছি (এগুলিই
``পর্যায়ক্রমিক ধারণা সম্পর্কে উপপাদ্য'')।
সবশেষে, প্রতিস্থাপন দিয়ে ট্রান্সফাইনাইট
পুনরাবৃত্তির উপপাদ্য প্রমাণ করেছি (এটাই
``সুভিত্তিক সম্পর্কের উপর আবেশী সংজ্ঞা'')।

এগুলি নিঃসন্দেহে কাম্য ফল। কিন্তু এই
ফলগুলি কি প্রতিস্থাপনকে \emph{ন্যায্য}
প্রতিপন্ন করার পক্ষে যথেষ্ট? \emph{না}।
অন্তত সরাসরি যথেষ্ট নয়।

একটি সরল সমস্যা দেখি। প্রতিস্থাপনের কিছু
কাম্য ফল দেখিয়েছি বটে, কিন্তু সেগুলির
অনেকগুলি অন্য পথেও পাওয়া যেত। বিষয়টি
যতটা জানা উচিত, ততটা পরিচিত নয়;
তাই একটু থেমে অবস্থা ব্যাখ্যা করি।

স্তরতত্ত্ব বা সংক্ষেপে $\LT$ নামে একটি
সরল সেটতত্ত্ব আছে।\footnote{$\LT$-র প্রথম
রূপগুলি \citet{Montague1965} ও
\citet{Scott1974} দিয়েছেন; \citet{Potter2004}
তত্ত্বটিকে সরল করে বইয়ের আকারে আলোচনা
করেছেন; সম্প্রতি \citet{ButtonLT1} $\LT$-কে
আরও সরল করেছেন।} $\LT$-র স্বতঃসিদ্ধ
কেবল বিস্তৃতিত্ব, পৃথকীকরণ এবং এই দাবি যে
প্রত্যেক সেট কোনো একটি \emph{স্তরের} উপসেট;
``স্তর'' এমনভাবে সংজ্ঞায়িত যে সেগুলি আমাদের
পরিচিত $V_\alpha$-গুলির মতো আচরণ করে।
তাই $\ZF$ থেকে $\LT$ অনুসৃত হয়; কিন্তু
$\LT$, $\ZF$-এর তুলনায় \emph{অনেক}
দুর্বল। বস্তুত, $\LT$ থেকে জোড়া, ঘাতসেট,
অসীমতা বা প্রতিস্থাপন—কোনোটিই পাওয়া যায় না।
$\LT$-র সঙ্গে অসীমতা ও ঘাতসেট যোগ করে
যে তত্ত্ব পাই তাকে $\Zr$ বলি। এতে জোড়াও
পাওয়া যায়; ফলে $\Zr$ অন্তত $\Z$-এর
সমান শক্তিশালী। কিন্তু আসলে $\Zr$,
$\Z$-এর চেয়ে কঠোরভাবে শক্তিশালী:
এটি আরও দাবি করে যে প্রত্যেক সেটের
স্তরাঙ্ক আছে (এই জন্যই $\Zr$ নামের
প্রস্তাব)। বস্তুত, $\Zr$ থেকে পাওয়া যায়:
ক্রমসংখ্যার সম্পূর্ণ সন্তোষজনক তত্ত্ব;
সুক্রমবিন্যস্ত পর্যায়ে স্তরক্রমের বিভাজনের ফল;
$\in$-আবেশের প্রমাণ; এবং ট্রান্সফাইনাইট
পুনরাবৃত্তির একটি \emph{রূপ}।

সুতরাং বুলস এটি না জানলেও, তাঁর বলা
সব কাম্য ফলই প্রতিস্থাপন \emph{ছাড়া}
পাওয়া যেত; শুধু $\Z$-এর বদলে
$\Zr$ ব্যবহার করতে হতো।

(এ সব সত্ত্বেও প্রচলিত পথ ধরে কেন
$\LT$ ও $\Zr$-এর বদলে $\ZF$
শেখালাম? দুটি কারণ। প্রথমত, নিছক
ঐতিহাসিক কারণে $\LT$ দিয়ে শুরু করা
বরং অপ্রচলিত; সেটতত্ত্বের আরও প্রচলিত
আলোচনা পড়ার উপযোগী করতে চেয়েছি।
দ্বিতীয়ত, $\LT$ ও $\Zr$ বোঝার
প্রস্তুতি হলে সরাসরি \citealt{Potter2004}
ও \citealt{ButtonLT1} পড়তে পারবে।)

অবশ্যই, $\Zr$ যেহেতু $\ZF$-এর তুলনায়
কঠোরভাবে দুর্বল, তাই কিছু ফল $\ZF$
প্রমাণ করলেও $\Zr$ মীমাংসা করে না।
সেই ধরনের কোনো ফল দেখিয়ে প্রতিস্থাপনের
বাহ্যিক ন্যায্যতা প্রতিষ্ঠার চেষ্টা করা যায়।
কিন্তু $\Zr$ ব্যবহার করতে জানলে এমন
উদাহরণ পাওয়া বেশ কঠিন, যা
(a) প্রতিস্থাপন দিয়ে মীমাংসিত হয়,
অন্যভাবে নয়, এবং (b) স্বজ্ঞায় সত্য।
(এ বিষয়ে আরও দেখো
\citealt[\S13.2]{Potter2004}।)

মূল কথা এই: প্রতিস্থাপনের পক্ষে জোরালো
বাহ্যিক ন্যায্যতা দিতে এমন একটি ফল
দরকার, যা প্রতিস্থাপন \emph{ছাড়া}
পাওয়া \emph{যায় না}। সে কাজ সহজ নয়।

প্রতিস্থাপনের \emph{শুধু} বাহ্যিক ন্যায্যতা
দেওয়ার যেকোনো চেষ্টায় আরেকটি সমস্যা
দেখা দেয়। সেটি হয়তো প্রথমটির চেয়েও
মৌলিক। বুলস শুধু বলেননি যে প্রতিস্থাপনের
অনেক কাম্য ফল আছে; তিনি আরও বলেছেন,
এর ``(দৃশ্যত) কোনো অকাম্য'' ফল নেই।
কিন্তু ``দৃশ্যত'' সতর্কতাটি নিশ্চিতভাবেই
অত্যন্ত গুরুত্বপূর্ণ।

এখানে পৌঁছানোর পথটি মনে করি: অবাধ
উপলব্ধি অসংগতিতে পড়েছিল; তার জবাবে
আমরা সেটের সঞ্চয়ী-পর্যায়ক্রমিক ধারণা
গ্রহণ করি। ধারণাটির সঙ্গে এমন একটি
আখ্যান আছে, যা আশা করি তার সঙ্গতি
সম্পর্কে আমাদের আশ্বস্ত করে। কিন্তু সেই
আখ্যানের \emph{ভিতর} থেকে প্রতিস্থাপনের
ন্যায্যতা দিতে না পারলে, $\ZF$ সঙ্গতিপূর্ণ
বলে বিশ্বাস করার (এখনও) কোনো কারণ নেই।
অথবা বলা ভালো: \emph{এখনও} কেউ কোনো
বিরোধ খুঁজে পাননি—এই (হয়তো নিছক
ঘটনাপ্রসূত) তথ্য ছাড়া $\ZF$-র সঙ্গতি
বিশ্বাসের কারণ নেই। ঠিক এই অর্থেই
বুলসের মন্তব্যটি দাঁড়ায়: ``(দৃশ্যত)
$\ZF$ সঙ্গতিপূর্ণ''। এর চেয়ে দৃঢ়
সঙ্গতি-আশ্বাস দাবি করা উচিত।

প্রতিস্থাপনের \emph{সম্পূর্ণ} বাহ্যিক
ন্যায্যতার যেকোনো চেষ্টাই এই সমস্যায়
পড়বে, অর্থাৎ যে যুক্তি কেবল $\ZF$-এর
(জানা) ফলের উপর নির্ভর করে। তাই
প্রতিস্থাপনের \emph{অন্তর্নিহিত} ন্যায্যতা
খুঁজব: এমন যুক্তি, যা দেখায় যে সেট
সম্বন্ধে আমাদের আখ্যান যেন ``আগেই''
প্রতিস্থাপনে আমাদের অঙ্গীকারবদ্ধ করেছে।

\end{document}
